\chapter{A complete two-trait worked example}

\section{Write the operator in the eigenbasis of \(\G\)}

Choose coordinates aligned with the eigenvectors of \(\G\), so
\begin{equation}
\G=
\begin{bmatrix}
g_1&0\\0&g_2
\end{bmatrix},
\qquad g_1\ge g_2>0.
\end{equation}
Let
\begin{equation}
\Hh=
\begin{bmatrix}
h_{11}&h_{12}\\h_{12}&h_{22}
\end{bmatrix}.
\end{equation}
Then
\begin{equation}
\G^{1/2}=
\begin{bmatrix}
\sqrt{g_1}&0\\0&\sqrt{g_2}
\end{bmatrix}.
\end{equation}
Multiply in two steps:
\begin{align}
\G^{1/2}\Hh
&=
\begin{bmatrix}
\sqrt{g_1}h_{11}&\sqrt{g_1}h_{12}\\
\sqrt{g_2}h_{12}&\sqrt{g_2}h_{22}
\end{bmatrix},\\[0.5em]
\Sop=\G^{1/2}\Hh\G^{1/2}
&=
\begin{bmatrix}
g_1h_{11}&\sqrt{g_1g_2}h_{12}\\
\sqrt{g_1g_2}h_{12}&g_2h_{22}
\end{bmatrix}.
\label{eq:S2D}
\end{align}

\inwords The diagonal entries combine variance and curvature along the genetic axes. The off-diagonal entry is curvature that mixes the two genetic axes, scaled by their geometric-mean variance.

\section{Derive the eigenvector angle}

Write a generic symmetric two-by-two matrix as
\begin{equation}
\Sop=
\begin{bmatrix}
a&b\\b&c
\end{bmatrix}.
\end{equation}
A unit vector at angle \(\phi\) is
\begin{equation}
\mathbf v=
\begin{bmatrix}\cos\phi\\\sin\phi\end{bmatrix}.
\end{equation}
The Rayleigh quotient is
\begin{align}
r(\phi)
&=\mathbf v^\top\Sop\mathbf v\\
&=a\cos^2\phi+2b\sin\phi\cos\phi+c\sin^2\phi.
\end{align}
Use
\begin{equation}
\cos^2\phi=\frac{1+\cos2\phi}{2},
\quad
\sin^2\phi=\frac{1-\cos2\phi}{2},
\quad
2\sin\phi\cos\phi=\sin2\phi,
\end{equation}
to obtain
\begin{equation}
r(\phi)=\frac{a+c}{2}+\frac{a-c}{2}\cos2\phi+b\sin2\phi.
\end{equation}
Differentiate:
\begin{equation}
\frac{\dd r}{\dd\phi}=-(a-c)\sin2\phi+2b\cos2\phi.
\end{equation}
At an eigenvector direction, the Rayleigh quotient is stationary, so
\begin{equation}
-(a-c)\sin2\phi+2b\cos2\phi=0.
\end{equation}
Rearranging,
\begin{equation}
\tan2\phi=\frac{2b}{a-c}.
\end{equation}
The quadrant-safe expression is
\begin{equation}
\boxed{\phi=\half\operatorname{atan2}(2b,a-c).}
\label{eq:2DeigenAngle}
\end{equation}
Substituting the entries from \cref{eq:S2D},
\begin{equation}
\boxed{
\tan2\phi=
\frac{2\sqrt{g_1g_2}\,h_{12}}
{g_1h_{11}-g_2h_{22}}.}
\label{eq:2DchannelAngle}
\end{equation}

\begin{mathteach}
The numerator is rotational forcing: it mixes the two axes. The denominator is the separation between their uncoupled response rates. Rotation is large when mixing is strong or when the two rates are nearly tied.
\end{mathteach}

\begin{geombox}
The two-dimensional formula already contains the later perturbation theory. The off-diagonal term tries to turn the axes; the eigenvalue gap resists that turn. Near a tie, a small off-diagonal perturbation can cause a large angle.
\end{geombox}

\section{Map the whitened angle back to phenotype space}

The whitened eigenvector is
\begin{equation}
\mathbf v=
\begin{bmatrix}\cos\phi\\\sin\phi\end{bmatrix}.
\end{equation}
The phenotype-space response direction is
\begin{equation}
\mathbf u=\G^{1/2}\mathbf v
=
\begin{bmatrix}
\sqrt{g_1}\cos\phi\\
\sqrt{g_2}\sin\phi
\end{bmatrix}.
\end{equation}
If \(\alpha\) is its ordinary Euclidean angle,
\begin{equation}
\tan\alpha
=\frac{u_2}{u_1}
=\sqrt{\frac{g_2}{g_1}}\tan\phi.
\label{eq:phenotypeAngle2D}
\end{equation}

\inwords Whitening and phenotype space do not assign the same Euclidean angle when genetic variance is anisotropic. The phenotype direction is stretched more along high-variance axes when mapped back from whitened coordinates.

\section{Three informative special cases}

\subsection{Isotropic curvature}
If \(\Hh=h\Ii\), then \(h_{12}=0\), \(h_{11}=h_{22}=h\), and \cref{eq:2DchannelAngle} gives \(\phi=0\) when \(g_1>g_2\). The leading channel is \(\mathbf g_{\max}\).

\subsection{Aligned but anisotropic curvature}
If \(h_{12}=0\), the matrices share axes. There is no smooth rotation. The leading direction is axis 1 when
\begin{equation}
g_1h_{11}>g_2h_{22}
\end{equation}
and axis 2 when the inequality reverses. At
\begin{equation}
g_1h_{11}=g_2h_{22},
\end{equation}
the eigenvalues are equal and the rank-one leading direction is undefined.

\begin{biobox}
Selection can change the strength of response modes without rotating them. Rotation requires a change in trait combinations, represented locally by cross-channel terms. A population can therefore experience strong selection and major eigenvalue change while its channel orientation remains constant.
\end{biobox}

\subsection{Misaligned curvature}
If \(h_{12}\neq0\), curvature couples the genetic axes and the leading channel rotates smoothly. The rotation becomes most sensitive as \(g_1h_{11}-g_2h_{22}\) approaches zero.

\begin{figure}[H]
\centering
\includegraphics[width=0.72\textwidth]{figures/figS1_channel_rotation_checked.pdf}
\caption{Two-trait illustration of rotation away from \(\mathbf g_{\max}\) as curvature anisotropy increases. Genetic correlation changes the orientation and susceptibility of the joint operator.}
\label{fig:2Drotation}
\end{figure}
